
\section{Introduction}\label{sec:introduction}

Insects play a critical role in ecosystems as pollinators, decomposers, and indicators of biodiversity.
However, insect populations are declining globally, and monitoring them at scale is difficult.
Passive acoustic recording is a practical approach: many insect species produce species-specific sounds that can be captured in the field.
Automated recognition of these sounds enables large-scale, cost-effective biodiversity monitoring.

The task of classifying insect species from audio is challenging.
Species within the same genus produce acoustically similar calls, recordings vary in quality, and datasets are often imbalanced.
InsectSet66~\cite{insectset66zenodo} provides a curated benchmark of 66 Orthoptera and Cicadidae species with official train, validation, and test splits.

Fai\ss{} and Stowell~\cite{faiss2023adaptive} proposed a convolutional neural network (CNN) trained on log-mel spectrograms as the reference system for this dataset.
Their 5-layer mel-CNN achieves 82\% test accuracy with data augmentation (colored noise and impulse responses).
Classical machine learning methods, while generally less accurate than deep models, offer interpretability and are useful for understanding which acoustic features drive classification.

In this project we:
(1) replicate the log-mel CNN baseline from Fai\ss{} and Stowell~\cite{faiss2023adaptive};
(2) extract MFCC features and apply k-means clustering to visualize the acoustic feature space;
(3) train SVM and Random Forest classifiers on MFCC statistics;
(4) analyze model interpretability through RF feature importance and confusion patterns.
